\section{Introduction}\label{sec:intro}

Iterative proof refinement---whether carried out by a human mathematician,
a formal verification system, or a large language model---proceeds by
composing a sequence of transformations on a proof state.  Each
transformation may introduce a small error, and understanding how these
errors accumulate is essential for predicting whether a multi-step reasoning
process will converge.  Despite the practical importance of this question
in automated theorem proving \cite{gottesman1996}, no mathematical framework
currently characterizes the relationship between the algebraic structure of
error composition and the resulting convergence rate.

The study of error propagation on algebraic structures has a rich history
outside of proof theory.  In robotics and navigation, Barrau and Bonnabel
\cite{barrau2014} showed that systems evolving on matrix Lie groups
satisfying the \emph{group affine property} enjoy a remarkable
\emph{log-linear} error dynamics: the Lie logarithm of the nonlinear error
satisfies a linear differential equation, valid for arbitrarily large
errors.  Wang and Chirikjian \cite{wang2008} developed second-order error
propagation formulas on non-abelian Lie groups, revealing that
non-commutativity introduces coupling between mean and covariance evolution
that is absent on abelian groups.  Li et al.\ \cite{li2023} extended this
framework to stochastic affine group systems with external disturbances.

In quantum information theory, Gottesman \cite{gottesman1996} proved that
the stabilizer group of any quantum error-correcting code must be abelian,
providing a concrete example where abelian subgroup structure is necessary
for systematic error control.  Arvind et al.\ \cite{arvind2002} constructed
error-correcting codes from abelian subgroups of the non-abelian error
group, and Zhao and Miyake \cite{zhao2023} used group symmetries for error
mitigation.  These results suggest a general principle: abelian structure
enables better error management.

On the probabilistic side, the theory of random walks on groups, developed
extensively by Diaconis and collaborators \cite{diaconis1988}, connects the
spectral gap of a transition operator to mixing times via the group's
representation theory.  For abelian groups, all irreducible representations
are one-dimensional, yielding a complete eigenvalue decomposition.

In this paper, we bring these threads together by modeling proof state
transformations as elements of $\GL_d(\R)$ and analyzing error accumulation
through the lens of Lyapunov exponents and spectral gaps.  Our main
contributions are as follows.

\begin{enumerate}[leftmargin=*,itemsep=2pt,topsep=4pt]
\item We prove a \emph{Lyapunov exponent ordering theorem}
  (Theorem~\ref{thm:lyapunov}): abelian (simultaneously diagonalizable)
  error operators have strictly lower Lyapunov exponents than non-abelian
  ones, with the gap quantified by the expected squared commutator norm
  at order $\Theta(\eps^2)$.

\item We establish a \emph{spectral gap comparison}
  (Theorem~\ref{thm:spectral_gap}) for lazy random walks on finite
  groups, showing that the dihedral group $D_n$ has spectral gap
  asymptotically four times that of the cyclic group $\Z_{2n}$.

\item We prove that the \emph{Diaconis upper bound} is tighter for abelian
  groups by a factor equal to the maximal irreducible representation
  dimension (Theorem~\ref{thm:diaconis}).

\item We apply these results to iterative proof refinement
  (Theorem~\ref{thm:convergence}), identifying distinct robustness
  thresholds for abelian and non-abelian error structures.
\end{enumerate}

The paper is organized as follows.  Section~\ref{sec:prelim} introduces
notation and background results.  Section~\ref{sec:results} contains the
statements and proofs of all main results.  Section~\ref{sec:discussion}
discusses implications, connections to prior work, and open problems.
Section~\ref{sec:conclusion} summarizes our contributions.
